\section{Related Work}

\subsection{PDP, PoR, and Public Auditing}

PDP and PoR provide the foundation for remote data integrity verification. Ateniese et al. introduced PDP, allowing a client to verify outsourced data possession without downloading the whole file~\cite{ateniese2007pdp}. Juels and Kaliski proposed PoR, which emphasizes recoverability rather than mere possession~\cite{juels2007por}. Shacham and Waters further reduced proof size and enabled public verification~\cite{shacham2013por}. Dynamic PDP extends this line of work to authenticated insertions, deletions, and modifications~\cite{erway2015dpdp}.

These protocols answer how to verify remote data at low communication cost. They do not, by themselves, provide a complete audit governance layer. In data circulation, the audit request, response time, verification result, recovery action, and responsibility assignment must also be recorded in a tamper-evident way. TrustAuditFlow therefore treats PDP/PoR as candidate off-chain proof modules and focuses on how the generated evidence is anchored, confirmed, and connected with recovery.

\subsection{Blockchain-Based Auditing and Smart Contracts}

Blockchain-based auditing uses an append-only ledger to constrain the audit process. Xue et al. studied blockchain-assisted identity-based public auditing against malicious auditors~\cite{xue2019ibpa}. Zhang et al. used unpredictable on-chain information to generate challenges and record auditor behavior against procrastination~\cite{zhang2021cpvpa}. Dredas implements challenge generation, deposits, penalties, and audit records through smart contracts~\cite{fan2020dredas}. CBPIV uses a consortium blockchain and PBFT consensus to record public integrity verification for IoT cloud storage~\cite{lin2022cbpiv}.

These systems support the general pattern of off-chain proof generation and on-chain evidence anchoring. However, when audit events are frequent, committing every challenge and every fine-grained result to the chain may create avoidable throughput pressure. In addition, many systems stop at pass/fail integrity checking and do not fully connect audit failures with recovery, accountability, and future node selection.

\subsection{Erasure Coding and Distributed Storage Reliability}

An integrity audit that only detects corruption but cannot trigger recovery has limited operational value. Erasure coding provides fault tolerance with lower redundancy than naive replication~\cite{balaji2018erasure}. BFT-Store combines Reed-Solomon coding with BFT/PBFT assumptions to reduce the storage burden of full replication in permissioned blockchain storage~\cite{qi2021bftstore}. PartitionChain explores reliable storage partitioning, aggregation, and reputation mechanisms for permissioned chains~\cite{du2023partitionchain}. These systems motivate the recovery component of TrustAuditFlow: audit failures should update not only the verification state, but also the recoverability state.

\subsection{Consortium Consensus and Trusted Computing}

Data circulation often involves known institutions and explicit access control, making consortium chains more suitable than open proof-of-work chains for audit records. PBFT is a classical Byzantine fault-tolerant protocol for state-machine replication~\cite{castro2002pbft}. HotStuff improves view-change complexity and responsiveness~\cite{yin2019hotstuff}, while DAG-BFT protocols separate data dissemination from ordering for high-throughput settings~\cite{danezis2022narwhal,spiegelman2022bullshark}. BLoP combines blockchain, trusted computing, and privacy-preserving computation, and uses PBFT-style consensus as a lower-energy alternative to PoW in consortium scenarios~\cite{blop2024trusted}. TrustAuditFlow builds on that full-node model and adds a risk-aware lightweight path rather than replacing the BLoP baseline.
